\chapter{전자기파와 맥스웰 작용}

\lectureoverview{먼저 소스가 없는 맥스웰 방정식을 풀어 평면파를 구하고, 그 방정식으로부터 파의 편광과 장의 기하학적 형태, 전파 속도를 결정한다. 이어 더 근본적인 질문을 살펴본다. 어떤 국소적이고 로렌츠 불변이며 게이지 불변인 작용이 동역학적 맥스웰 방정식을 만들어 내는가? 그 답은 전자기장, 전하 보존, 게이지 불변성을 하나의 구조로 엮어 준다.}

\section{맥스웰 방정식의 두 계열}

평면 전자기파에서 시작하자. 여기서 \emph{plane}은 기하학적 용어이지, 꾸밈이 없다는 뜻의 \emph{plain}이 아니다. 먼저 구체적인 질문을 던져 보자. 이런 파는 어떤 모양이며, 그 형태는 맥스웰 방정식에서 어떻게 나오는가?

파를 다루기 전에 방정식을 두 열로 정리해 두면 편리하다. 진공에서는 다음과 같다.
\begin{align}
\nabla\cdot\mathbf B&=0,
&
\dot{\mathbf B}&=\nabla\times\mathbf E,
\label{eq:homogeneous-vector}
\\
\nabla\cdot\mathbf E&=0,
&
\dot{\mathbf E}&=-\,\nabla\times\mathbf B.
\label{eq:dynamical-vector-vacuum}
\end{align}
부호가 대부분 양수인 계량과
\begin{equation}
\eta_{\mu\nu}=\operatorname{diag}(-1,1,1,1)
\end{equation}
강의에서 사용하는 규약
\begin{equation}
F_{\mu\nu}
=
\partial_\mu A_\nu-\partial_\nu A_\mu,
\qquad
F_{0i}=E_i,
\qquad
F_{ij}=\epsilon_{ijk}B_k.
\label{eq:course-f-convention}
\end{equation}
을 적용하면, 같은 두 계열을 다음과 같이 쓸 수 있다.
\begin{align}
\partial_\mu F_{\nu\sigma}
+\partial_\nu F_{\sigma\mu}
+\partial_\sigma F_{\mu\nu}
&=0,
\label{eq:bianchi-l9}
\\
\partial_\mu F^{\mu\nu}
&=0.
\label{eq:vacuum-maxwell-covariant}
\end{align}

\begin{figure}[H]
\centering
\includegraphics[width=\linewidth]{lecture_09_figure_01.png}
\caption{진공 맥스웰 방정식의 두 계열과 그 공변식. 순환 방정식은 항등식이고, 축약된 방정식은 동역학을 기술한다.}
\par\smallskip
\textit{강의 화면: 00:03:30.}
\end{figure}

두 행의 차이는 중요하다. 방정식 \eqref{eq:bianchi-l9}은 퍼텐셜 \(A_\mu\)이 존재한다는 사실에서 항등적으로 따라온다. 각 항이 이계 미분이고 항들이 쌍을 이루어 상쇄되기 때문이다. 따라서 이는 비앙키 항등식이다. 반면 방정식 \eqref{eq:vacuum-maxwell-covariant}은 운동 방정식이다. 전자기 라그랑지안을 바꾸면 달라지는 것은 바로 이 계열이다.

칠판에 적힌 작은 관찰 하나가 이 장에서 두 가지 역할을 한다. 맥스웰 방정식에 나오는 모든 장에는 미분이 작용한다. 이 사실은 평면파의 상대 위상을 제약하고, 뒤에서는 게이지 불변성을 자연스럽게 만들어 준다.

\section{\texorpdfstring{\(z\)}{z}축을 따라 진행하는 평면파}

좌표축을 선택하여 파동이 \(z\)축을 따라 진행하도록 하자.
이렇게 해도 일반성을 잃지 않는다. 다른 방향으로 진행하는 파동도 축을 회전하면 이 형태로 나타낼 수 있다. 파수를 \(k\)라 하고,
\begin{equation}
\lambda=\frac{2\pi}{k},
\end{equation}
진동수를 \(\omega\)라 하자. 운동 방정식을 적용하기 전에, 여섯 성분 모두에 공통 위상
\begin{equation}
\theta=kz-\omega t:
\end{equation}
을 사용하여 다음과 같이 쓰자.
\begin{align}
E_x(z,t)&=\epsilon_x\sin\theta,
&
B_x(z,t)&=\beta_x\sin\theta,
\nonumber\\
E_y(z,t)&=\epsilon_y\sin\theta,
&
B_y(z,t)&=\beta_y\sin\theta,
\label{eq:plane-wave-ansatz}\\
E_z(z,t)&=\epsilon_z\sin\theta,
&
B_z(z,t)&=\beta_z\sin\theta.
\nonumber
\end{align}
여섯 그리스 문자는 상수 진폭이다.

\begin{figure}[H]
\centering
\includegraphics[width=\linewidth]{lecture_09_figure_02.png}
\caption{맥스웰 방정식이 허용되지 않는 진폭을 제거하기 전의 여섯 성분 평면파 가정.}
\par\smallskip
\textit{강의 화면: 00:08:45.}
\end{figure}

왜 \(\mathbf E\)와 \(\mathbf B\)에 같은 사인 의존성을 먼저 가정하는가? 한 장은 \(\sin\theta\)에 비례하고 다른 장은 \(\cos\theta\)에 비례한다고 해 보자. 컬 방정식의 한쪽에서 시간 미분을 하고 다른 쪽에서 공간 미분을 하면 서로 다른 삼각함수가 나온다. 그러면 모든 \(z\)와 \(t\)에 대해 등식이 성립할 수 없다. 일정한 공통 위상은 허용되지만, \(z\)의 원점을 옮기면 이를 없앨 수 있다. 따라서 이 가정은 별도의 동역학 법칙이 아니다. 미분 방정식이 요구하는 위상 관계를 가장 간단히 나타낸 것이다.

아직 맥스웰 방정식은 하나도 사용하지 않았다. 진행 방향과 파장, 편리한 위상 원점만 정했다. 이제 방정식이 어떤 진폭이 남을지 결정한다.

\subsection{발산 제약 조건}

장들이 \(z\)과 \(t\)에만 의존하므로,
\begin{equation}
\nabla\cdot\mathbf E
=
\partial_xE_x+\partial_yE_y+\partial_zE_z
=
\partial_zE_z.
\end{equation}
진공에서는 모든 위상에서 이 값이 0이어야 하므로,
\begin{equation}
\epsilon_z=0.
\label{eq:epsilon-z-zero}
\end{equation}
자기장에도 같은 논리를 적용하면
\begin{equation}
\beta_z=0.
\label{eq:beta-z-zero}
\end{equation}
두 장 모두 진행 방향 성분이 없다. 파동은 횡파이다.

횡방향 \(xy\)평면에는 회전 자유도가 남아 있다. \((\epsilon_x,\epsilon_y)\)는 보통의 2차원 벡터이므로 횡축을 회전하여
\begin{equation}
\epsilon_y=0,
\qquad
\epsilon_x\ne0.
\label{eq:polarization-axis}
\end{equation}
로 둘 수 있다. 이는 선형 편광을 선택한 것이지 물리학을 제한한 것이 아니다.

\begin{figure}[H]
\centering
\includegraphics[width=\linewidth]{lecture_09_figure_03.png}
\caption{종방향 진폭이 사라지고 횡방향 전기장 축을 선택했다.}
\par\smallskip
\textit{강의 화면: 00:14:30.}
\end{figure}

\subsection{컬 방정식}

\(\dot{\mathbf E}=-\nabla\times\mathbf B\)의 \(x\)성분을 보자. 장은 \(z\)에만 의존하므로,
\begin{equation}
\dot E_x
=
-\left(\nabla\times\mathbf B\right)_x
=
\partial_zB_y.
\end{equation}
파동 가정을 대입하면
\begin{align}
-\omega\epsilon_x\cos\theta
&=
k\beta_y\cos\theta,
\nonumber\\
\beta_y
&=
-\,\epsilon_x\frac{\omega}{k}.
\label{eq:beta-y-epsilon}
\end{align}
같은 맥스웰 방정식의 \(y\)성분은 다음과 같다.
\begin{equation}
\dot E_y
=
-\left(\partial_zB_x\right).
\end{equation}
편광을 선택했으므로 좌변은 0이다. 따라서
\begin{equation}
\beta_x=0.
\label{eq:beta-x-zero}
\end{equation}

\begin{figure}[H]
\centering
\includegraphics[width=\linewidth]{lecture_09_figure_04.png}
\caption{성분별 컬 계산으로 \(\epsilon_x\)와 \(\beta_y\)의 관계를 구하고 남은 자기장 진폭을 없앤다.}
\par\smallskip
\textit{강의 화면: 00:20:15.}
\end{figure}

따라서 기하학적 배치는 다음과 같이 정해진다.
\begin{equation}
\mathbf E\parallel\hat{\mathbf x},
\qquad
\mathbf B\parallel\hat{\mathbf y},
\qquad
\mathbf k\parallel\hat{\mathbf z}.
\end{equation}
전기장, 자기장, 진행 방향은 서로 수직이다.

\begin{figure}[H]
\centering
\includegraphics[width=0.94\linewidth]{lecture_09_figure_05.png}
\caption{횡방향 전기장과 자기장의 진동이 \(z\)축을 따라 진행하는 칠판 그림.}
\par\smallskip
\textit{강의 화면: 00:23:15.}
\end{figure}

다른 컬 방정식에서도 독립적인 진폭 관계가 나온다. 그 \(y\)성분은
\begin{equation}
\dot B_y=(\nabla\times\mathbf E)_y=\partial_zE_x,
\end{equation}
이므로
\begin{align}
-\omega\beta_y\cos\theta
&=
k\epsilon_x\cos\theta,
\nonumber\\
\epsilon_x
&=
-\,\beta_y\frac{\omega}{k}.
\label{eq:epsilon-beta-y}
\end{align}

\begin{figure}[H]
\centering
\includegraphics[width=\linewidth]{lecture_09_figure_06.png}
\caption{횡파 그림 위에 적힌 두 진폭 사이의 상호 관계.}
\par\smallskip
\textit{강의 화면: 00:25:00.}
\end{figure}

방정식 \eqref{eq:beta-y-epsilon}과 \eqref{eq:epsilon-beta-y}을 결합하면 분산 관계를 얻는다.
\begin{equation}
\left(\frac{\omega}{k}\right)^2=1.
\end{equation}
선택한 방향으로 진행하는 양의 진동수 파에 대해서는
\begin{equation}
\omega=k
\qquad
(c=1).
\label{eq:omega-k}
\end{equation}
따라서 공통 위상은
\begin{equation}
\theta=k(z-t).
\end{equation}
가 된다. 통상적인 단위를 복원하면
\begin{equation}
\omega=ck,
\qquad
\theta=k(z-ct).
\label{eq:restore-c-wave}
\end{equation}
즉 파형은 빛의 속도로 변형 없이 이동한다. 칠판에서 사용하는 단위계에서는 전기장과 자기장의 진폭 크기가 같다. 두 진폭 사이의 상대적인 마이너스 부호는 방향 규약이다. 이는 선택한 \(y\)축이 어느 쪽을 가리키는지를 나타내며, 한 장이 다른 장보다 사분의 일 주기 늦다는 뜻은 아니다.

\begin{figure}[H]
\centering
\includegraphics[width=\linewidth]{lecture_09_figure_07.png}
\caption{\(\omega=k\)를 적용한 뒤 얻은 최종 평면파 형태.}
\par\smallskip
\textit{강의 화면: 00:26:45.}
\end{figure}

\begin{classroomqa}
\audiencequestion{\(\omega\)와 \(k\)는 단위가 서로 다르지 않나요?}
\lecturerresponse{통상적인 단위계에서는 그렇습니다. \(\omega\)는 시간의 역수이고 \(k\)는 길이의 역수입니다. \(c=1\)에서는 시간과 길이의 단위가 같습니다. \(c\)를 복원하면 \(\omega=k\)가 \(\omega=ck\)로 바뀌며, 이는 차원 분석의 요구와 정확히 일치합니다. \lecturetimestamp{00:25:55}}
\end{classroomqa}

\begin{classroomqa}
\audiencequestion{평면파에서 전기장과 자기장은 언제나 같은 위상인가요?}
\lecturerresponse{두 장은 같은 위상 의존성을 갖습니다. 부호를 포함한 진폭은 마이너스 부호만큼 다를 수 있는데, 이는 \(180^\circ\)의 반전이라고 표현할 수 있습니다. 그러나 맥스웰 방정식은 이 선형 편광 평면파에서 \(90^\circ\)만큼의 사인--코사인 위상차를 허용하지 않습니다. 원편광은 횡방향 성분들을 조합하여 만들기 때문에 더 세심한 표현이 필요하지만, 성분별로 적용되는 이 맥스웰 관계는 달라지지 않습니다. \lecturetimestamp{00:27:31}}
\end{classroomqa}

\begin{classroomqa}
\audiencequestion{두 장이 동시에 0을 지나면 에너지는 여전히 보존되나요?}
\lecturerresponse{그렇습니다. 파가 지나갈 때 고정된 한 지점의 에너지 밀도는 진동하지만, 에너지가 사라진 것이 아니라 이동한 것입니다. 한 시점에 한 영역에 있던 에너지는 잠시 뒤 그 옆 영역에 놓입니다. 포인팅 플럭스로 더 완전하게 표현되는 국소적 에너지 수송을 고려하면, 에너지 밀도의 변화와 보존 법칙은 양립합니다.
\lecturetimestamp{00:28:15}}
\end{classroomqa}

맥스웰 방정식은 선형이다. 진폭 \(\epsilon_x\)는 클 수도 작을 수도 있고, \(\beta_y\)는 그 진폭으로부터 정해지며, 해를 더하거나 빼도 다시 해가 된다. 따라서 서로 다른 방향으로 진행하는 파들을 중첩할 수 있다. 이 짧은 계산은 해 볼 가치가 있다. 전자기학을 설명하면서 전자기파를 근거 없이 인용된 결과로 남겨 두어서는 안 된다.

\begin{classroomqa}
\audiencequestion{처음부터 \(\mathbf E\)와 \(\mathbf B\)의 위상이 같다고 가정하고 유도한 건가요?}
\lecturerresponse{공통 위상을 가정 형태로 적기는 했지만, 그것은 독립적인 물리적 가정이 아닙니다. 컬 방정식의 양쪽에서 각각 장을 미분하므로, 방정식이 자명하지 않은 상대 위상을 배제합니다. 그 뒤 공통 위상 이동은 원점을 옮겨 없앨 수 있습니다. \lecturetimestamp{00:29:56}}
\end{classroomqa}

\begin{classroomqa}
\audiencequestion{두 번째 관계 \(\epsilon_x=-\beta_y\omega/k\)는 어디서 나왔나요?}
\lecturerresponse{첫 번째 관계는 \(\dot{\mathbf E}=-\nabla\times\mathbf B\)에서 나왔습니다. 두 번째는 같은 성분 계산에서 \(\mathbf E\)와 \(\mathbf B\)을 맞바꾸어 적용한 짝이 되는 방정식 \(\dot{\mathbf B}=\nabla\times\mathbf E\)에서 얻습니다. \(\omega^2=k^2\)를 구하려면 두 관계가 모두 필요합니다.
\lecturetimestamp{00:30:25}}
\end{classroomqa}

\section{미분 방정식만으로는 충분하지 않은 이유}

이제 특정 해에서 이론의 구조로 논의를 옮겨 가자. 최소 작용 원리는 에너지 보존과 운동량 보존, 그리고 이들이 시공간 대칭과 맺는 관계를 정리해 준다. 보존된다고 부를 만한 에너지도, 쓸모 있는 대칭 원리도 없는 미분 방정식은 얼마든지 만들 수 있다. 하지만 맥스웰 방정식은 라그랑지안에서 나와야 한다.

동역학을 기술하는 계열만 유도하면 된다.
\begin{equation}
\partial_\mu F^{\mu\nu}=J^\nu.
\label{eq:target-sourced-maxwell}
\end{equation}
비앙키 계열은 이미 \(F_{\mu\nu}\)의 정의에서 따라온다. 3차원 벡터 표기에서 우리의 규약은 다음을 준다.
\begin{align}
\nabla\cdot\mathbf E&=\rho,
\label{eq:gauss-source}
\\
\dot{\mathbf E}+\nabla\times\mathbf B&=-\,\mathbf J.
\label{eq:ampere-course-sign}
\end{align}
\eqref{eq:ampere-course-sign}의 마이너스 부호가 중요하다. 이는 \eqref{eq:target-sourced-maxwell}을 \eqref{eq:course-f-convention} 규약으로 썼을 때의 성분 형태이며, 강의 끝부분의 실시간 부호 확인에서 다시 확인하게 된다.

\begin{figure}[H]
\centering
\includegraphics[width=\linewidth]{lecture_09_figure_08.png}
\caption{파동 해에서 작용 원리로 넘어갈 때의 맥스웰 방정식 전체가 적힌 칠판.}
\par\smallskip
\textit{강의 화면: 00:32:45.}
\end{figure}

\section{국소적 상대론적 장 이론}

첫 번째 원리는 국소성이다. 작용은 같은 시공간 점에서의 장과 그 일차 미분으로 만든 밀도를 시공간 전체에 걸쳐 적분한 것이다.
\begin{equation}
S=\int d^4x\,\mathcal L(\phi,\phi_{,\mu}),
\qquad
\phi_{,\mu}\equiv\partial_\mu\phi.
\label{eq:local-action}
\end{equation}
쉼표는 미분을 간결하게 표기한 것이다. 두 번째 원리는 로렌츠 불변성이다. \(\mathcal L\)는 스칼라여야 한다.

\subsection{스칼라장으로 미리 연습하기}

4-벡터를 변분하기 전에 스칼라장 하나로 방법을 연습하자. 국소적인 로렌츠 스칼라는 퍼텐셜 \(U(\phi)\)와 두 미분의 수축에 의존할 수 있다.
\begin{equation}
\mathcal L_\phi
=
-\frac12\,\partial_\mu\phi\,\partial^\mu\phi-U(\phi).
\label{eq:scalar-lagrangian}
\end{equation}
부호가 대부분 양수인 계량을 사용하면
\begin{equation}
\partial_\mu\phi\,\partial^\mu\phi
=
-\dot\phi^{\,2}+|\nabla\phi|^2,
\end{equation}
이므로
\begin{equation}
\mathcal L_\phi
=
\frac12\dot\phi^{\,2}
-\frac12|\nabla\phi|^2
-U(\phi).
\end{equation}
시간 미분의 위 첨자는 아래 첨자일 때와 부호가 다르지만, 공간 미분의 위 첨자와 아래 첨자는 부호가 같다.

\begin{classroomqa}
\audiencequestion{미분 지수를 위에 두는 것과 아래에 두는 것은 어떤 의미가 있나요?}
\lecturerresponse{그렇습니다. 지수를 올리면 역계량과 수축합니다. 계량 부호 규약이 \((-+++)\)일 때,
\(\partial^0=-\partial_0\)인 반면
\(\partial^i=\partial_i\)입니다. 그래서 시간 미분 항의 부호는 세 공간 미분 항의 부호와 반대입니다.
\lecturetimestamp{00:38:02}}
\end{classroomqa}

장에 대한 오일러--라그랑주 방정식은 다음과 같다.
\begin{equation}
\partial_\mu\left(
\frac{\partial\mathcal L}{\partial\phi_{,\mu}}
\right)
=
\frac{\partial\mathcal L}{\partial\phi}.
\label{eq:field-el}
\end{equation}
\eqref{eq:scalar-lagrangian}에 대해서는
\begin{equation}
\frac{\partial\mathcal L_\phi}{\partial\phi_{,\mu}}
=
-\partial^\mu\phi,
\qquad
\frac{\partial\mathcal L_\phi}{\partial\phi}
=
-U'(\phi),
\end{equation}
따라서
\begin{equation}
\Box\phi=U'(\phi),
\qquad
\Box\equiv\partial_\mu\partial^\mu.
\end{equation}
동등하게 쓰면
\begin{equation}
\ddot\phi-\nabla^2\phi=-U'(\phi).
\label{eq:scalar-wave}
\end{equation}
칠판에는 \(x\) 미분항만 적혀 있고, \(y\)와 \(z\) 항은 생략되어 있다. 결과는 퍼텐셜로 인한 힘 항이 있는 익숙한 파동 방정식이다.

\begin{figure}[H]
\centering
\includegraphics[width=\linewidth]{lecture_09_figure_09.png}
\caption{국소성, 로렌츠 불변 스칼라 라그랑지안, 장 이론의 오일러--라그랑주 방정식.}
\par\smallskip
\textit{강의 화면: 00:41:00.}
\end{figure}

전자기장 계산도 같은 변분 절차를 따르지만, 장이 네 개이고 지수가 더 많다.

\section{게이지 불변성이 맥스웰 스칼라를 선택한다}

세 번째 원리는 게이지 불변성이다. 우선 \(J^\mu=0\)로 두고 자유장을 생각하자. 게이지 변환은 다음과 같다.
\begin{equation}
A_\mu\longrightarrow A_\mu+\partial_\mu\chi.
\label{eq:gauge-transform}
\end{equation}
혼합 이계 미분은 서로 교환되므로 장 텐서는 변하지 않는다.
\begin{align}
F_{\mu\nu}
&\longrightarrow
F_{\mu\nu}
+\partial_\mu\partial_\nu\chi
-\partial_\nu\partial_\mu\chi
\nonumber\\
&=F_{\mu\nu}.
\end{align}

로렌츠 불변성만으로는 스칼라 \begin{equation}
A_\mu A^\mu.
\end{equation}도 허용된다. 하지만 \eqref{eq:gauge-transform}에 따라 변하므로 맥스웰 이론에서는 허용되지 않는다. 이 식의 지수는 장 성분을 나타내며 미분을 뜻하지 않는다. 미분은 \(A_{\nu,\mu}=\partial_\mu A_\nu\)처럼 쉼표로 표시한다.

\begin{figure}[H]
\centering
\includegraphics[width=\linewidth]{lecture_09_figure_10.png}
\caption{게이지 변환과 게이지 불변이 아니어서 배제되는 로렌츠 스칼라 \(A_\mu A^\mu\).}
\par\smallskip
\textit{강의 화면: 00:43:40.}
\end{figure}

따라서 자유 라그랑지안은 \(F_{\mu\nu}\)로 구성해야 한다.
선형 대각합은 성립하지 않는다:
\begin{equation}
F^\mu{}_\mu=0.
\end{equation}
반대칭성 때문에 모든 대각 성분은 0이다:
\begin{equation}
F_{00}=F_{11}=F_{22}=F_{33}=0.
\end{equation}
지표를 올리면 부호가 바뀔 수는 있지만, 0이 다른 값으로 바뀌지는 않는다.

\begin{figure}[H]\centering
\includegraphics[width=\linewidth]{lecture_09_figure_11.png}\caption{반대칭 장 텐서의 선형 대각합을 시도한 결과와 사라지는 대각 성분.}
\par\smallskip
\textit{강의 화면: 00:47:00.}
\end{figure}

처음으로 0이 아닌 로런츠 스칼라는 이차항이다:
\begin{equation}
F_{\mu\nu}F^{\mu\nu}.
\label{eq:f2-scalar}
\end{equation}
수축 결과를 바로 인용하는 대신 풀어 써 보자. 혼합 성분에는 전기장이 들어 있다. 시간 지표를 올리면 마이너스 부호가 생기고, \(0i\)와 \(i0\) 두 순서가 모두 나타난다:
\begin{equation}
\sum_i\left(F_{0i}F^{0i}+F_{i0}F^{i0}\right)
=-2E^2.
\end{equation}
순수 공간 성분에는 자기장이 들어 있다. 공간 지표를 올려도 부호는 바뀌지 않으며, 여기서도 두 순서가 모두 나타난다:
\begin{equation}
\sum_{i<j}\left(F_{ij}F^{ij}+F_{ji}F^{ji}\right)
=2B^2.
\end{equation}
따라서
\begin{equation}
F_{\mu\nu}F^{\mu\nu}=2(B^2-E^2),
\qquad
-F_{\mu\nu}F^{\mu\nu}=2(E^2-B^2).
\label{eq:f2-eb}
\end{equation}

\begin{figure}[H]
\centering
\includegraphics[width=\linewidth]{lecture_09_figure_12.png}
\caption{전기장과 자기장의 크기로 나타낸 수축의 최종 결과.}
\par\smallskip
\textit{강의 화면: 00:51:00.}
\end{figure}

\begin{classroomqa}\audiencequestion{\(F_{00}\)도 0인가?}\lecturerresponse{그렇다. \(F_{00}\)를 포함해 반대칭 텐서의 모든 대각 성분은 0이다. \(\mathbf E\)을 나타내는 비대각 혼합 성분과 \(\mathbf B\)를 나타내는 비대각 공간 성분만 기여한다.
\lecturetimestamp{00:51:36}}\end{classroomqa}

관례적인 정규화는 다음과 같다.
\begin{equation}
\boxed{
\mathcal L_{\mathrm{EM}}
=
-\frac14F_{\mu\nu}F^{\mu\nu}
}
\label{eq:maxwell-lagrangian}
\end{equation}
또는
\begin{equation}
\mathcal L_{\mathrm{EM}}
=
\frac12(E^2-B^2).
\label{eq:maxwell-lagrangian-eb}
\end{equation}
마이너스 부호는 전기장 항에 익숙한 양의 부호를 준다. 계수 \(1/4\)은 반대칭 성분 쌍이 두 번 나타나는 것을 보정해 관례적인 계수 \(1/2\)을 남긴다. 자유 작용 전체에 0이 아닌 상수를 곱해도 소스가 없는 방정식은 바뀌지 않지만, 이 정규화는 표준 소스 결합에 맞는다.

\section{자유 전자기장의 변분}

독립 변수인 장은 다음과 같다.
\begin{equation}
A_0,\qquad A_x,\qquad A_y,\qquad A_z.
\end{equation}
각 성분 \(A_\nu\)에 대해 오일러--라그랑주 방정식은 다음과 같다.
\begin{equation}
\partial_\mu\left(
\frac{\partial\mathcal L}{\partial A_{\nu,\mu}}
\right)
=
\frac{\partial\mathcal L}{\partial A_\nu},
\qquad
A_{\nu,\mu}\equiv\partial_\mu A_\nu.
\label{eq:vector-el-comma}
\end{equation}
쉼표 표기에서 지수의 위치는 천천히 살펴볼 만하다. 첫 번째 지수는 \(A\)의 어떤 성분이 장인지 나타내고, 두 번째 지수는 미분하는 좌표를 나타낸다.

이 표기법에서
\begin{equation}
F_{\mu\nu}
=
A_{\nu,\mu}-A_{\mu,\nu}.
\label{eq:f-comma-consistent}
\end{equation}
이는 정확히 \eqref{eq:course-f-convention}이다. 쉼표로 표시한 지수의 순서를 말로 거꾸로 쓰면 중간에 마이너스 부호가 바뀐다. 성분별 계산을 해 보면 유용한 이유가 여기에 있다.

\begin{figure}[H]
\centering
\includegraphics[width=\linewidth]{lecture_09_figure_13.png}
\caption{네 퍼텐셜 성분을 독립적인 장으로 놓은 뒤 쉼표 표기로 나타낸 장 텐서.}
\par\smallskip
\textit{강의 화면: 00:55:30.}
\end{figure}

\begin{figure}[H]
\centering
\includegraphics[width=\linewidth]{lecture_09_figure_14.png}
\caption{벡터 퍼텐셜 미분으로 명시적으로 쓴 맥스웰 라그랑지안.}
\par\smallskip
\textit{강의 화면: 00:56:45.}
\end{figure}

\subsection{대표 미분 하나}

미분항 \(A_{x,y}=\partial_yA_x\)를 하나 고르자. 전체 합에서 \(xy\)항과 \(yx\)항의 기여는 다음과 같다.
\begin{align}
\mathcal L_{\mathrm{EM}}
&\supset
-\frac14\left(F_{xy}F^{xy}+F_{yx}F^{yx}\right)
\nonumber\\
&=
-\frac12F_{xy}^2
\nonumber\\
&=
-\frac12\left(A_{x,y}-A_{y,x}\right)^2.
\label{eq:isolated-xy-piece}
\end{align}
제곱하면 부호가 사라지므로 마지막 등식은 문제가 없다.
\(A_{x,y}-A_{y,x}=F_{yx}=-F_{xy}\).

\begin{figure}[H]
\centering
\includegraphics[width=\linewidth]{lecture_09_figure_15.png}
\caption{대표 미분항 \(A_{x,y}\)와 분리해 낸 반대칭 성분의 제곱.}
\par\smallskip
\textit{강의 화면: 01:00:45.}
\end{figure}

미분하면
\begin{align}
\frac{\partial\mathcal L_{\mathrm{EM}}}
{\partial A_{x,y}}
&=
-\left(A_{x,y}-A_{y,x}\right)
\nonumber\\
&=
F_{xy}.
\label{eq:component-variation}
\end{align}
칠판에서는 미분 지수가 칠판의 순서를 따르는 텐서에 마이너스 부호를 붙여 이 결과를 요약한다. 강의 전체에서 사용하는 정의 \eqref{eq:course-f-convention}에 따르면 불변인 식은 다음과 같다.
\begin{equation}
\frac{\partial\mathcal L_{\mathrm{EM}}}
{\partial(\partial_\mu A_\nu)}
=
-F^{\mu\nu}.
\label{eq:maxwell-canonical-derivative}
\end{equation}
예를 들어 \(\partial_yA_x\)의 경우에는 \(-F^{yx}=F^{xy}\)가 되어, 정확히 \eqref{eq:component-variation}와 일치한다. 따라서 겉으로 보이는 부호 차이는 지수 순서를 정리하는 과정에서 생긴 것일 뿐이다.

\begin{figure}[H]
\centering
\includegraphics[width=\linewidth]{lecture_09_figure_16.png}
\caption{퍼텐셜 미분항 하나에 대해 미분하기 전 선택한 제곱을 전개한 모습.}
\par\smallskip
\textit{강의 화면: 01:03:00.}
\end{figure}

\begin{classroomqa}
\audiencequestion{성분 미분에서 \(1/2\) 인자는 어떻게 된 건가요?}
\lecturerresponse{제곱을 미분하면 2가 나오므로 \(1/2\)와 상쇄됩니다.
\[
\frac{d}{dx}\left[-\frac12x^2\right]=-x.
\]
앞서 \(1/4\)에서 \(1/2\)로 바뀐 것은 순서가 있는 두 텐서 성분 \(xy\)와 \(yx\)에서 비롯됐습니다. \lecturetimestamp{01:03:29}}
\end{classroomqa}

진공 라그랑지안에는 미분되지 않은 퍼텐셜이 없다.
\begin{equation}
\frac{\partial\mathcal L_{\mathrm{EM}}}{\partial A_\nu}=0.
\end{equation}
이를 \eqref{eq:vector-el-comma}에 대입하면
\begin{equation}
-\partial_\mu F^{\mu\nu}=0,
\end{equation}
또는 동등하게
\begin{equation}
\boxed{\partial_\mu F^{\mu\nu}=0.}
\label{eq:free-maxwell-from-action}
\end{equation}

\begin{figure}[H]
\centering
\includegraphics[width=\linewidth]{lecture_09_figure_17.png}
\caption{일반화한 성분별 결과를 오일러--라그랑주 방정식에 대입한다.}
\par\smallskip
\textit{강의 화면: 01:09:00.}
\end{figure}

작용 전체를 변분하면 같은 결론을 더 간결하게 얻는다.
\begin{align}
\delta S_{\mathrm{EM}}
&=
-\frac12\int d^4x\,
F^{\mu\nu}
\left(\partial_\mu\delta A_\nu
-\partial_\nu\delta A_\mu\right)
\nonumber\\
&=
-\int d^4x\,F^{\mu\nu}\partial_\mu\delta A_\nu
\nonumber\\
&=
\int d^4x\,
\left(\partial_\mu F^{\mu\nu}\right)\delta A_\nu,
\label{eq:global-maxwell-variation}
\end{align}
여기서는 경계항을 버렸다. 첫째 줄의 두 항은 반대칭성으로 합쳐진다. \(\delta A_\nu\)가 임의이므로 \eqref{eq:free-maxwell-from-action}가 따라온다.

\begin{classroomqa}
\audiencequestion{제곱항 하나만 떼어 낸 것은 자기장이 \(z\) 방향만 향한다고 가정한 건가요?}
\lecturerresponse{아닙니다. 살펴보기 위해 전체 합에서 한 항을 뽑아낸 것이며, 그 뒤에는 생략 부호가 들어가야 합니다. 나머지 공간 성분 쌍에서 다른 자기장 성분들이 나오고, 시간--공간 성분 쌍에서는 전기장 성분들이 나옵니다. \lecturetimestamp{01:08:05}}
\end{classroomqa}

\begin{classroomqa}
\audiencequestion{공간 지수 두 개 대신 시간 지수 하나를 선택하면 결과가 달라지나요?}
\lecturerresponse{일반적인 미분 결과는 같으며, 계량 부호는 지수를 올리는 과정에서 반영됩니다. 공간 지수 두 개를 선택하면 자기장 성분이 나오고, 시간 지수 하나와 공간 지수 하나를 선택하면 전기장 성분이 나옵니다. 반대칭 공식 \eqref{eq:maxwell-canonical-derivative}은 두 경우를 모두 포괄합니다.
\lecturetimestamp{01:08:46}}
\end{classroomqa}

\begin{classroomqa}
\audiencequestion{여전히 \(J^\mu=0\)라고 가정하고 있나요?}
\lecturerresponse{자유장 유도에서만 그렇습니다. 게이지 불변 운동항을 확인하려고 전류를 잠시 0으로 두었습니다. 이제 전류를 다시 넣고, 맥스웰 방정식의 우변을 만들어 내는 추가 항이 무엇인지 살펴보겠습니다. \lecturetimestamp{01:09:41}}
\end{classroomqa}

대수 자체는 어렵지 않지만, 그저 지켜보는 데 그치지 말고 직접 전개해야 한다. 지수의 연쇄를 한 번 따라가 보는 것만으로는 오래 기억에 남지 않지만, 손으로 직접 전개하면 남는다.

\section{보존되는 4-전류}

소스는 4-전류다.
\begin{equation}
J^\mu=(\rho,\mathbf J).
\end{equation}
시간 성분은 전하 밀도이고, 공간 성분은 서로 수직인 세 면을 통과하는 전하 흐름이다. 또는 크기가 작은 고정 상자를 생각하면
\begin{equation}
\frac{\partial\rho}{\partial t}
+\nabla\cdot\mathbf J=0.
\label{eq:continuity-vector-l9}
\end{equation}
상자 안의 전하가 줄어들 수 있는 것은 경계를 통해 빠져나갈 때뿐이다. 공변 표기로는
\begin{equation}
\boxed{\partial_\mu J^\mu=0.}
\label{eq:continuity-covariant-l9}
\end{equation}

\begin{figure}[H]
\centering
\includegraphics[width=\linewidth]{lecture_09_figure_18.png}
\caption{소스 결합을 도입하기 전에 다시 살펴보는 전하 밀도, 공간 전류, 작은 상자 해석.}
\par\smallskip
\textit{강의 화면: 01:14:45.}
\end{figure}

\section{소스 결합과 게이지 불변성}

전류와 퍼텐셜에 각각 선형인 가장 간단한 국소 로렌츠 스칼라는 다음과 같다.
\begin{equation}
S_{\mathrm{int}}
=
-\int d^4x\,J^\mu A_\mu.
\label{eq:interaction-action}
\end{equation}
전체 부호는 규약에 따른다. 강의에서는 전하 부호의 역사적 선택을 결국 벤저민 프랭클린 탓으로 돌릴 수도 있다는 농담을 한다. 핵심은 일관된 부호 규약 하나를 어디서나 사용해야 한다는 점이다.

처음에는 \eqref{eq:interaction-action}가 게이지 불변으로 보이지 않는다. \eqref{eq:gauge-transform}에 따라
\begin{equation}
\delta_\chi S_{\mathrm{int}}
=
-\int d^4x\,J^\mu\partial_\mu\chi.
\label{eq:gauge-variation-interaction}
\end{equation}

\begin{figure}[H]
\centering
\includegraphics[width=\linewidth]{lecture_09_figure_19.png}
\caption{게이지 변환 \(A_\mu\mapsto A_\mu+\partial_\mu\chi\)에 따른 상호작용 작용과 그 변화.}
\par\smallskip
\textit{강의 화면: 01:16:45.}
\end{figure}

기본적인 방식으로 작동 원리를 보려면 수축을 전개하면 된다.
\begin{equation}
J^\mu\partial_\mu\chi
=
J^0\partial_0\chi
+J^1\partial_1\chi
+J^2\partial_2\chi
+J^3\partial_3\chi.
\end{equation}
전류가 큰 실험실 안에 국한되어 있거나 충분히 빠르게 감소하여 경계에서 사라진다고 가정하자. 각 항을 부분 적분한다. 예를 들어 첫 번째 공간 좌표에 대해서는
\begin{align}
-\int d^4x\,J^1\partial_1\chi
&=
-\int dt\,dx^2\,dx^3
\left[J^1\chi\right]_{x^1=-\infty}^{x^1=+\infty}
\nonumber\\
&\quad
+\int d^4x\,\chi\,\partial_1J^1
\nonumber\\
&=
\int d^4x\,\chi\,\partial_1J^1.
\end{align}
시간 항과 나머지 두 공간 항에도 같은 연산을 적용한다.

\begin{figure}[H]
\centering
\includegraphics[width=\linewidth]{lecture_09_figure_20.png}
\caption{부분 적분하기 전에 전개한 네 개의 게이지 변분 항.}
\par\smallskip
\textit{강의 화면: 01:19:15.}
\end{figure}

\begin{figure}[H]
\centering
\includegraphics[width=\linewidth]{lecture_09_figure_21.png}
\caption{게이지 함수에서 전류로 미분을 옮기는 대표적인 부분 적분.}
\par\smallskip
\textit{강의 화면: 01:20:45.}
\end{figure}

네 성분을 모두 모으면
\begin{equation}
\delta_\chi S_{\mathrm{int}}
=
\int d^4x\,\chi\,\partial_\mu J^\mu.
\label{eq:gauge-variation-continuity}
\end{equation}
따라서
\begin{equation}
\partial_\mu J^\mu=0
\quad\Longrightarrow\quad
\delta_\chi S_{\mathrm{int}}=0.
\end{equation}
가정한 경계 조건 아래에서, 소스 상호작용이 게이지 불변일 필요충분조건은 전류가 국소적으로 보존된다는 것이다.

\begin{classroomqa}\audiencequestion{\(\nabla\cdot\mathbf J=0\)만 요구하면 충분한가?}\lecturerresponse{아니다. 네 항이 모두 필요하다:
\[
\partial_\mu J^\mu
=
\dot\rho+\nabla\cdot\mathbf J=0.
\]
첫 번째 항은 고정된 상자 안에서 전하 밀도가 변하는 양이다. 그 변화로 보상된다면 공간 발산이 0이 아니어도 된다. \lecturetimestamp{01:22:00}}\end{classroomqa}

\section{소스가 있는 맥스웰 방정식}

전자기 라그랑지안 밀도 전체는 다음과 같다.
\begin{equation}
\boxed{
\mathcal L
=
-\frac14F_{\mu\nu}F^{\mu\nu}
-J^\mu A_\mu.
}
\label{eq:full-maxwell-lagrangian}
\end{equation}
소스 항에는 \(A_\nu\)의 미분이 없으므로 \eqref{eq:maxwell-canonical-derivative}는 바꾸지 않고, 다음을 기여한다.
\begin{equation}
\frac{\partial\mathcal L}{\partial A_\nu}=-J^\nu.
\end{equation}
따라서 오일러--라그랑주 방정식은
\begin{equation}
-\partial_\mu F^{\mu\nu}=-J^\nu,
\end{equation}
또는
\begin{equation}
\boxed{\partial_\mu F^{\mu\nu}=J^\nu.}
\label{eq:sourced-maxwell-from-action}
\end{equation}

\begin{figure}[H]
\centering
\includegraphics[width=\linewidth]{lecture_09_figure_22.png}
\caption{대표적인 오일러--라그랑주 미분과 함께 표시한 자유장 항과 소스 항.}
\par\smallskip
\textit{강의 화면: 01:25:00.}
\end{figure}

보존되지 않는 전류를 주어진 소스로 사용하면 두 가지 문제가 동시에 생긴다. 상호작용 작용은 게이지 불변이 아니게 되고, 소스가 있는 맥스웰 방정식은 서로 모순된다. 경계를 통과하는 보상 흐름이 없는데 전하가 영역 안에서 그냥 사라질 수는 없다.

\section{맥스웰 방정식은 보존 법칙을 되돌려 준다}

장 방정식에서 호환성을 바로 확인할 수 있다. \eqref{eq:sourced-maxwell-from-action}에 \(\partial_\nu\)를 작용시키면
\begin{equation}
\partial_\nu J^\nu
=
\partial_\nu\partial_\mu F^{\mu\nu}
=0,
\label{eq:antisymmetry-implies-continuity}
\end{equation}
미분의 두 지수 \(\mu\leftrightarrow\nu\)에 대해 대칭이고 \(F^{\mu\nu}\)는 반대칭이기 때문이다.

강의에서는 이를 삼차원 벡터 형태로도 확인하는데, 그 과정에서 아직 남아 있는 부호 오류가 드러난다. 가우스 법칙에서 시작하자:
\begin{equation}
\rho=\nabla\cdot\mathbf E.
\end{equation}
시간에 대해 미분하면
\begin{equation}
\dot\rho=\nabla\cdot\dot{\mathbf E}.
\label{eq:rho-dot-div-edot}
\end{equation}
칠판에 잠정적으로 적힌 방정식은 다음과 같다.
\begin{equation}
\dot{\mathbf E}
=
-\nabla\times\mathbf B+\mathbf J.
\label{eq:provisional-wrong-source-sign}
\end{equation}
이 식의 발산을 취하면 다음을 얻는다.
\begin{equation}
\dot\rho=\nabla\cdot\mathbf J,
\end{equation}
회전의 발산은 사라지기 때문이다. 이는 \eqref{eq:continuity-vector-l9}와 모순된다.

\begin{figure}[H]
\centering
\includegraphics[width=\linewidth]{lecture_09_figure_23.png}
\caption{잘못된 연속 방정식 부호를 내는 잠정적인 플러스 \(\mathbf J\) 항 방정식.}
\par\smallskip
\textit{강의 화면: 01:27:00.}
\end{figure}

불일치는 숨겨져 있지 않다. 계산을 실시간으로 되짚고, 필요한 연속 방정식을 밝힌 뒤, 청중 한 사람이 방정식을 확인하여 소스 항의 부호를 찾아낸다. 규약에 일관된 성분 방정식은
\begin{equation}
\boxed{
\dot{\mathbf E}+\nabla\times\mathbf B=-\,\mathbf J,
}
\label{eq:corrected-source-sign}
\end{equation}
또는
\begin{equation}
\dot{\mathbf E}
=
-\nabla\times\mathbf B-\mathbf J.
\end{equation}
이다. 이제 방정식 \eqref{eq:rho-dot-div-edot}과 \eqref{eq:corrected-source-sign}를 결합하면
\begin{align}
\dot\rho
&=
\nabla\cdot
\left(-\nabla\times\mathbf B-\mathbf J\right)
\nonumber\\
&=
-\nabla\cdot\mathbf J,
\end{align}
을 얻으며, 이는 정확히 국소 전하 보존이다.

\begin{figure}[H]
\centering
\includegraphics[width=\linewidth]{lecture_09_figure_24.png}
\caption{수정된 소스 부호와 그로부터 되찾은 연속 방정식.}
\par\smallskip
\textit{강의 화면: 01:31:00.}
\end{figure}

\begin{classroomqa}
\audiencequestion{우변의 전류 항은 \(-\mathbf J\)여야 하지 않나요?}
\lecturerresponse{강의 전체에서 사용한 장 텐서 규약에 따르면 \(\partial_\mu F^{\mu\nu}=J^\nu\)의 공간 성분은
\(\dot{\mathbf E}+\nabla\times\mathbf B=-\mathbf J\)입니다.
이렇게 고치면 발산 계산에서 \(\dot\rho+\nabla\cdot\mathbf J=0\)가 나오며, 앞서 공변 작용을 유도한 결과는 달라지지 않습니다. \lecturetimestamp{01:30:17}}
\end{classroomqa}

\section{어떤 원리가 근본적인가?}

\begin{classroomqa}\audiencequestion{논리가 순환적인가? 작용을 게이지 불변으로 만들기 위해 전하 보존을 사용한 다음, 작용으로 맥스웰 방정식을 유도했고, 맥스웰 방정식은 전하 보존을 함의한다.}\lecturerresponse{서로 양립하는 진술이지, 결론을 전제로 삼는 순환 논증은 아니다. 관측된 맥스웰 방정식에서 출발해 연속 방정식을 발견한 다음, \(-J^\mu A_\mu\)가 게이지 불변 결합임을 알아볼 수 있다. 또는 국소성, 로런츠 불변성, 게이지 불변성과 작용에서 출발할 수도 있다. 게이지 불변성은 보존 전류를 요구하고, 작용은 맥스웰 방정식을 준다. 또는 전하가 국소적으로 사라지지 않는다는 물리적 사실을 출발점으로 삼을 수도 있다. 어느 경로를 택해도 같은 구조에 도달한다. \lecturetimestamp{01:31:11}}\end{classroomqa}

가장 간결한 구성 경로를 묻는다면, 국소성, 로런츠 불변성, 게이지 불변성과 라그랑지안 원리의 결합은
\begin{equation}
-\frac14F_{\mu\nu}F^{\mu\nu}-J^\mu A_\mu
\end{equation}
를 선택하며, 따라서 맥스웰의 동역학 방정식으로 이어진다. 이 삼각형의 어느 요소가 더 근본적인지 말하는 것은 어느 정도 취향의 문제다. 중요한 것은 세 요소가 서로 맞물린다는 사실이다.

\section{‘\emph{게이지}’라는 말의 유래}

\begin{classroomqa}\audiencequestion{수학적 내용은 분명한데, 여기서 \emph{게이지}라는 말은 무슨 뜻인가?}\lecturerresponse{이 명칭은 시대착오적인 역사적 용어다. 게이지는 기준 압력에 대한 압력처럼, 선택된 영점을 가리킬 수 있다. 헤르만 바일은 처음에 이 개념을 국소적인 길이의 기준과 연결했지만, 그 해석은 올바른 물리가 아니었다. 자유도가 양자 파동함수의 위상과 올바르게 연결된 뒤에도 이 용어는 살아남았다. 지금의 고전적 계산에서는 \(A_\mu\)를 \(\partial_\mu\chi\)만큼 바꾸어도 기술상의 영점만 바뀌고 \(F_{\mu\nu}\)는 바뀌지 않는다. 더 깊은 양자적 의미는 뒤에서 다룬다.
\lecturetimestamp{01:34:49}}\end{classroomqa}

역사적 명칭 때문에 핵심 진술이 흐려져서는 안 된다. \eqref{eq:gauge-transform}로 서로 연결된 퍼텐셜은 같은 전자기장을 나타낸다. 따라서 게이지 자유도는 퍼텐셜을 사용한 기술에서의 중복성이며, 장 텐서는 게이지 불변이다.

\section{맥스웰 방정식의 두 계열은 왜 비대칭적으로 보이는가}

\begin{classroomqa}\audiencequestion{전기장과 자기장이 대칭적이라면, 왜 맥스웰 방정식의 한 계열은 비앙키 항등식으로 나타나고 다른 계열은 운동 방정식에서 나오는가?}\lecturerresponse{소스가 없는 전자기학에서는 쌍대 퍼텐셜을 도입하여 두 역할을 맞바꿀 수 있다. 그러면 운동 방정식이었던 것이 비앙키 항등식이 되고, 그 반대도 성립한다. 물리적 비대칭성은 소스에서 비롯된다. 전하는 관측되었지만 기본 자기 홀극은 관측되지 않았다. 전하를 서로 바꾸는 정확한 통상 대칭이 있다면, 그에 대응하는 스펙트럼을 가진 자기 전하도 있을 것으로 예상되지만 관측되지 않는다. 더 정교한 이론에는 여전히 자기 전하나 전기--자기 쌍대성이 포함될 수 있지만, 이는 지금의 유도 범위를 벗어난다. \lecturetimestamp{01:37:11}}\end{classroomqa}

따라서 소스가 없는 방정식에는 전기--자기 교환 구조가 있지만, 관측된 소스 부문은 전기적 기술을 택한다. 이는 단순한 표기 관례 이상이지만, 자유 평면파에서 보이는 대칭성을 무효화하지는 않는다.

\section{지금까지 세운 구조}

강의에서 따라온 전개를 이제 마무리할 수 있다:
\begin{enumerate}
\item 맥스웰 방정식은 평면파가 횡파이고 전기장과 자기장이 서로 직교하며 분산 관계가 \(\omega=ck\)임을 강제한다.
\item 비앙키 방정식은 퍼텐셜에서 따라 나오는 항등식이고, 수축된 방정식은 동역학적 방정식이다.
\item 국소성과 로런츠 불변성은 작용 밀도로 이끌며, 게이지 불변성은 \(A_\mu A^\mu\)을 배제하고 \(F_{\mu\nu}\)로 만든 스칼라를 선택한다.
\item 자유 작용
      \(-\tfrac14F_{\mu\nu}F^{\mu\nu}\)는 진공 맥스웰 방정식을 준다.
\item 상호작용 항 \(-J^\mu A_\mu\)는 전류가 연속 방정식을 만족할 때, 그리고 그때에만 게이지 불변이다.
\item 소스가 있는 작용은
      \(\partial_\mu F^{\mu\nu}=J^\nu\)를 주며, 반대칭성으로부터
      \(\partial_\mu J^\mu=0\)가 다시 나온다.
\end{enumerate}

이 안에는 라그랑주 역학, 상대론적 장 이론, 전기역학, 로런츠 불변성, 게이지 불변성, 국소 전하 보존이라는 많은 구조가 하나의 논증으로 엮여 있다. 이를 소화할 수 있는 믿을 만한 방법은 방정식부터 다시 시작해 각 미분과 부호를 손으로 하나씩 따라가는 것이다.